% GENERATED FILE. Edit canonical lessons, then run npm run generate:books.
% canonical-source-hash: fnv1a64:51ea15e6a64c6d1c
% canonical-lessons: GU-C311-samacar, GU-C311-chapu, GU-C311-mathalu, GU-C311-rediyo, GU-C311-janavvu

\chapter{The Morning News}
\label{ch:gu-the-morning-news}
\hlchaptermodality{311}

\begin{chapteropening}
\textbf{By the end of this chapter:} \emph{I can say what the news said this morning.}

Complete the last lesson of chapter 311: I can say what the news said this morning.
\end{chapteropening}

\section[\texorpdfstring{\gu{સમાચાર} (samācār)}{સમાચાર (samācār)}]{\gu{સમાચાર} (samācār) — the news}

\label{lesson:GU-C311-samacar}



\begin{quote}
Before the new one: say the Gujarati for to be at a loss, then the Gujarati for to disperse.
\end{quote}

\subsection*{You'll want to know: \gu{સમાચાર}}
\textbf{\gu{સમાચાર}} — \emph{samācār} — \textquotedblleft{}the news\textquotedblright{}.

\begin{quote}
\gu{મેં} \gu{સવારે} \gu{સમાચાર} \gu{સાંભળ્યા}.
\end{quote}

\begin{quote}
\emph{mẽ savāre samācār sā̃bhaḷyā.}
\end{quote}

\begin{quote}
I listened to the news this morning.
\end{quote}

A masculine noun, used in the plural.

\begin{cousinweb}
\textbf{\gu{સમાચાર}} (\emph{samācār}) is a tatsama, Sanskrit \emph{samācāra}, \textquotedblleft{}procedure, conduct, custom\textquotedblright{} (Monier-Williams); the sense \textquotedblleft{}news\textquotedblright{} is the modern language's. The verb with it is plural: \textbf{\gu{સમાચાર} \gu{સાંભળ્યા}}.
\end{cousinweb}

\begin{grammarlens}[title={What you've built}]
One more word for telling what happened.
\end{grammarlens}

\subsection*{Your turn}
\begin{itemize}
\raggedright
  \item \emph{Say it:} \emph{samācār}
  \item \emph{Say it:} \emph{samācār}, once more
  \item \emph{Say it:} \emph{mẽ savāre samācār sā̃bhaḷyā.} — I listened to the news this morning.
  \item \emph{Recall:} say the Gujarati for to be at a loss, then the Gujarati for to disperse, then say \emph{samācār} again
\end{itemize}

\begin{tcolorbox}[breakable,colback=teal!4,colframe=teal!35!black,title={Before you move on}]
What does \gu{સમાચાર} mean? (\textquotedblleft{}The news\textquotedblright{}.) Say it once more.
\end{tcolorbox}
\section[\texorpdfstring{\gu{છાપું} (chāpũ)}{છાપું (chāpũ)}]{\gu{છાપું} (chāpũ) — a newspaper}

\label{lesson:GU-C311-chapu}



\begin{quote}
Before the new one: say the Gujarati for to disperse, then the Gujarati for the news.
\end{quote}

\subsection*{You'll want to know: \gu{છાપું}}
\textbf{\gu{છાપું}} — \emph{chāpũ} — \textquotedblleft{}a newspaper\textquotedblright{}.

\begin{quote}
\gu{મારા} \gu{દાદા} \gu{રોજ} \gu{સવારે} \gu{છાપું} \gu{વાંચે} \gu{છે}.
\end{quote}

\begin{quote}
\emph{mārā dādā roj savāre chāpũ vā̃ce che.}
\end{quote}

\begin{quote}
My grandfather reads the newspaper every morning.
\end{quote}

A neuter noun.

\begin{cousinweb}
\textbf{\gu{છાપું}} (\emph{chāpũ}) is built on \textbf{\gu{છાપવું}}, to print, from \emph{*chapp}, \textquotedblleft{}press, cover\textquotedblright{}, a form CDIAL reconstructs (CDIAL 4994): the printed sheet. \textbf{\gu{છાપામાં}}, in the paper.
\end{cousinweb}

\begin{grammarlens}[title={What you've built}]
One more word for telling what happened.
\end{grammarlens}

\subsection*{Your turn}
\begin{itemize}
\raggedright
  \item \emph{Say it:} \emph{chāpũ}
  \item \emph{Say it:} \emph{chāpũ}, once more
  \item \emph{Say it:} \emph{mārā dādā roj savāre chāpũ vā̃ce che.} — My grandfather reads the newspaper every morning.
  \item \emph{Recall:} say the Gujarati for to disperse, then the Gujarati for the news, then say \emph{chāpũ} again
\end{itemize}

\begin{tcolorbox}[breakable,colback=teal!4,colframe=teal!35!black,title={Before you move on}]
What does \gu{છાપું} mean? (\textquotedblleft{}A newspaper\textquotedblright{}.) Say it once more.
\end{tcolorbox}
\section[\texorpdfstring{\gu{મથાળું} (mathāḷũ)}{મથાળું (mathāḷũ)}]{\gu{મથાળું} (mathāḷũ) — a headline}

\label{lesson:GU-C311-mathalu}



\begin{quote}
Before the new one: say the Gujarati for the news, then the Gujarati for a newspaper.
\end{quote}

\subsection*{You'll want to know: \gu{મથાળું}}
\textbf{\gu{મથાળું}} — \emph{mathāḷũ} — \textquotedblleft{}a headline\textquotedblright{}.

\begin{quote}
\gu{છાપાનું} \gu{મથાળું} \gu{હતું}: \gu{પોલીસે} \gu{ઠગને} \gu{ઝડપ્યો}.
\end{quote}

\begin{quote}
\emph{chāpānũ mathāḷũ hatũ: polīse ṭhagne jhaḍapyo.}
\end{quote}

\begin{quote}
The headline read: the police nab a swindler.
\end{quote}

A neuter noun.

\begin{cousinweb}
\textbf{\gu{મથાળું}} (\emph{mathāḷũ}) comes from \emph{*mastatalaka}, a form CDIAL reconstructs on Sanskrit \emph{mastaka}, the head, and CDIAL lists Gujarati \emph{mathāḷũ}, \textquotedblleft{}top, heading\textquotedblright{} (CDIAL 9927). In a newspaper it is the line at the head of a story.
\end{cousinweb}

\begin{grammarlens}[title={What you've built}]
One more word for telling what happened.
\end{grammarlens}

\subsection*{Your turn}
\begin{itemize}
\raggedright
  \item \emph{Say it:} \emph{mathāḷũ}
  \item \emph{Say it:} \emph{mathāḷũ}, once more
  \item \emph{Say it:} \emph{chāpānũ mathāḷũ hatũ: polīse ṭhagne jhaḍapyo.} — The headline read: the police nab a swindler.
  \item \emph{Recall:} say the Gujarati for the news, then the Gujarati for a newspaper, then say \emph{mathāḷũ} again
\end{itemize}

\begin{tcolorbox}[breakable,colback=teal!4,colframe=teal!35!black,title={Before you move on}]
What does \gu{મથાળું} mean? (\textquotedblleft{}A headline\textquotedblright{}.) Say it once more.
\end{tcolorbox}
\section[\texorpdfstring{\gu{રેડિયો} (reḍiyo)}{રેડિયો (reḍiyo)}]{\gu{રેડિયો} (reḍiyo) — a radio}

\label{lesson:GU-C311-rediyo}



\begin{quote}
Before the new one: say the Gujarati for a newspaper, then the Gujarati for a headline.
\end{quote}

\subsection*{You'll want to know: \gu{રેડિયો}}
\textbf{\gu{રેડિયો}} — \emph{reḍiyo} — \textquotedblleft{}a radio\textquotedblright{}.

\begin{quote}
\gu{રેડિયો} \gu{પર} \gu{સમાચાર} \gu{આવતા} \gu{હતા}.
\end{quote}

\begin{quote}
\emph{reḍiyo par samācār āvtā hatā.}
\end{quote}

\begin{quote}
The news was on the radio.
\end{quote}

A masculine noun.

\begin{cousinweb}
\textbf{\gu{રેડિયો}} (\emph{reḍiyo}) is borrowed from English \emph{radio}, and it is masculine: \textbf{\gu{રેડિયો} \gu{પર}}, on the radio.
\end{cousinweb}

\begin{grammarlens}[title={What you've built}]
One more word for telling what happened.
\end{grammarlens}

\subsection*{Your turn}
\begin{itemize}
\raggedright
  \item \emph{Say it:} \emph{reḍiyo}
  \item \emph{Say it:} \emph{reḍiyo}, once more
  \item \emph{Say it:} \emph{reḍiyo par samācār āvtā hatā.} — The news was on the radio.
  \item \emph{Recall:} say the Gujarati for a newspaper, then the Gujarati for a headline, then say \emph{reḍiyo} again
\end{itemize}

\begin{tcolorbox}[breakable,colback=teal!4,colframe=teal!35!black,title={Before you move on}]
What does \gu{રેડિયો} mean? (\textquotedblleft{}A radio\textquotedblright{}.) Say it once more.
\end{tcolorbox}
\section[\texorpdfstring{\gu{જણાવવું} (jaṇāvvũ)}{જણાવવું (jaṇāvvũ)}]{\gu{જણાવવું} (jaṇāvvũ) — to inform}

\label{lesson:GU-C311-janavvu}



\begin{quote}
Before the new one: say the Gujarati for a headline, then the Gujarati for a radio.
\end{quote}

\subsection*{You'll want to know: \gu{જણાવવું}}
\textbf{\gu{જણાવવું}} — \emph{jaṇāvvũ} — \textquotedblleft{}to inform\textquotedblright{}.

\begin{quote}
\gu{પોલીસે} \gu{લોકોને} \gu{જણાવ્યું} \gu{કે} \gu{રસ્તો} \gu{બંધ} \gu{છે}.
\end{quote}

\begin{quote}
\emph{polīse lokone jaṇāvyũ ke rasto bandh che.}
\end{quote}

\begin{quote}
The police told people that the road was closed.
\end{quote}

Past, with \gu{તેણે} (he or she …): \gu{તેણે} \gu{બધાને} \gu{જણાવ્યું}, \emph{teṇe badhāne jaṇāvyũ}.

\begin{cousinweb}
\textbf{\gu{જણાવવું}} (\emph{jaṇāvvũ}) belongs with Sanskrit \emph{jānāti}, \textquotedblleft{}knows\textquotedblright{}, where CDIAL lists the Gujarati verb, \textquotedblleft{}informs\textquotedblright{} (CDIAL 5193): to make someone know. It sits beside \textbf{\gu{જાણવું}}, to know.
\end{cousinweb}

\begin{grammarlens}[title={What you've built}]
One more word for telling what happened.
\end{grammarlens}

\subsection*{Your turn}
\begin{itemize}
\raggedright
  \item \emph{Say it:} \emph{jaṇāvvũ}
  \item \emph{Say it:} \emph{jaṇāvvũ}, once more
  \item \emph{Say it:} \emph{polīse lokone jaṇāvyũ ke rasto bandh che.} — The police told people that the road was closed.
  \item \emph{Recall:} say the Gujarati for a headline, then the Gujarati for a radio, then say \emph{jaṇāvvũ} again
\end{itemize}

\begin{tcolorbox}[breakable,colback=teal!4,colframe=teal!35!black,title={Before you move on}]
What does \gu{જણાવવું} mean? (\textquotedblleft{}To inform\textquotedblright{}.) Say it once more.
\end{tcolorbox}
